\section{Resonace, centralizers and cocharacter graphs}\label{section: cocharacter graphs and resonance}

We first introduce in \S \ref{section: resonance} the notion of \textit{resonance} inside $T \times \Gm^{\rot}$. Intuitively, two points in $T$ are resonant with respect to a scalar $\zeta \in \Gm^{\rot}$ if they multiplicatively differ by integer powers of $\zeta$. This gives rise to the \textit{resonant Weyl group} $W_{[x, \zeta]}$ and the associated \textit{resonant subgroup} $L_{[x, \zeta]}$. 
The notion of resonance appears when describing fixed-points of $(x, \zeta)$ in the affine Grassmannian and computing the corresponding fibers of its equivariant $K$-theory, as well as in the context of linkage principles in classical and quantum categories $\O$.

We then in \S \ref{subsection: centralizers in loop groups} describe the centralizer of $(x, \zeta)$ in $\lo G$ via its closed root subsystem, and discuss stabilizers of coordinate fixed points $t^{\lambda} \in \Gr_G$ inside it. We further consider the adjoint action of the resonant cocharacter $\omega$ on these objects, which we call \emph{slanting}. 
%
For simply connected $G$, we futher describe in \S \ref{subsection: weyl group stabilizers} the stabilizer of $(x, \zeta)$ inside the affine Weyl group.

We finally discuss in \S \ref{section: cocharacter graphs} the \textit{cocharacter graphs} $\Gamma$ inside $T \times \Gm^{\rot} \times T$ for a simply connected $G$. We explain how $\Gamma$ puts the resonant combinatorics into a family over $T \times \Gm^{\rot}$. Ultimately, $\Gamma$ turns out to be essential for our description of the whole $K$-theory ring of the affine Grassmannian and affine Schubert varieties inside it, because restriction to $\Gamma$ realizes equivariant localization to fixed points. We discuss the variant for reductive $G$ afterwards in \S \ref{section: cocharacter graphs for reductive groups}.

Altogether, the discussion in this section serves as a useful tool in the rest of the paper.








\subsection{Resonance and centralizers}\label{section: resonance}
To eventually describe fixed-points of elements $(x, \zeta) \in T \times \Gm^{\rot}$ on the affine Grassmannian, we introduce the notion of \emph{resonance}, \emph{resonant subgroups} and \emph{resonant cocharacter}.

\subsubsection{Resonance.}\label{subsubsection: resonance}
Let $(x, \zeta) \in T \times \Gm^{\rot}$. 
A point $x \in T$ is called \emph{$\zeta$-integral} if it is of the form $x=\lambda(\zeta)$ for some $\lambda \in \Z\Phi_{\bullet} \leq X_{\bullet}$. 
We say that $x, y \in T$ are \emph{$\zeta$-resonant} if there is some $\lambda \in \Z\Phi_{\bullet}$ such that $y = \lambda(\zeta) \cdot x$. In words, two points are $\zeta$-resonant if their ratio is $\zeta$-integral. 

We denote the equivalence relation given by $\zeta$-resonance as $x \sim_{\zeta} y$. For given $(x, \zeta) \in T \times \Gm^{\rot}$, we denote the equivalence class of $x$ under $\zeta$-resonance by $[x]_{\zeta} = [x, \zeta]$. We will use the latter notation to avoid too many subscripts.

The resonance relation is induced by the orbits of the \emph{$\zeta$-integral action} $\Z\Phi^{\bullet} \acts T$, $\lambda \cdot x = \lambda(\zeta) \cdot x$.

\subsubsection{Resonant subgroups.}
Recall the discussion of root subsystems from \S \ref{section: root subsystems and centralizers}.
Depending on $(x, \zeta) \in T \times \Gm^{\rot}$, consider the closed root subsystems
\begin{equation}
    \Phi^{\bullet}_x \subseteq \Phi^{\bullet}_{[x, \zeta]} \subseteq \Phi^{\bullet}_G
\end{equation}
given by
\begin{equation}\label{equation: resonant root system}
    \Phi^{\bullet}_x := \{\beta \in \Phi^{\bullet}_G \mid \beta(x)=1 \},
    \qquad \text{and} \qquad
    \Phi^{\bullet}_{[x, \zeta]} := \{\beta \in \Phi^{\bullet}_G \mid \exists m \in \Z, \ \beta(x)= \zeta^m \}.
\end{equation}
Indeed, both subsets $\Phi^{\bullet}_x$, $\Phi^{\bullet}_{[x, \zeta]}$, clearly satisfy (i), (ii) from Notation \ref{notation: closed root subsystems}.


We denote the corresponding regular subgroups by
\begin{equation}
L_x \leq L_{[x, \zeta]} \leq G.   
\end{equation}
We call them the \emph{stabilizer subgroup} and \emph{resonant subgroup}.
Their Weyl groups are denoted
\begin{equation}\label{equation: resonant weyl group}
    W_x \leq W_{[x, \zeta]} \leq W.
\end{equation}
We call them the \emph{stabilizer Weyl group} and \emph{resonant Weyl group}.
These are the reflection subgroups associated to the above closed root subsystems $\Phi^{\bullet}_{x}$ and $\Phi^{\bullet}_{[x, \zeta]}$. 


The resonant root subsystem $\Phi^{\bullet}_{[x, \zeta]}$, resonant subgroup $L_{[x, \zeta]}$ and resonant Weyl group $W_{[x, \zeta]}$ clearly depend, up to canonical isomorphism, only on the resonance class $[x, \zeta]$ of the element $(x, \zeta)$. 
%
The notation is thus consistent.
%
When the point $(x, \zeta)$ is clear from the context, we abbreviate
\begin{equation*}
L := L_{[x, \zeta]} \leq G, \qquad \qquad \Phi^{\bullet}_L := \Phi^{\bullet}_{[x, \zeta]} \subseteq \Phi^{\bullet}_G, \qquad \qquad W_L := W_{[x, \zeta]} \leq W. 
\end{equation*}



\subsubsection{Resonant cocharacter.}
Consider $(x, \zeta)$ as above.
Let $\beta \in \Phi^{\bullet}_L$ be a resonant root. By definition of $L$, there exists an integer $m_{\beta} \in \Z$ such that $\beta(x) = \zeta^{m_{\beta}}$. Put differently, we have $m_{\beta} \in \log_{\zeta}(\beta(x))$. The integer $m_{\beta}$ is unique modulo $\ell = \ord(\zeta)$. We may package such integers as follows.

\begin{construction}\label{notation: additive function m}
We choose an additive function   
\begin{equation*}
   \mm : \Phi^{\bullet}_L \to \Z, \qquad \beta \mapsto m_{\beta}
\end{equation*}
such that $\beta(x) = \zeta^{m_{\beta}}$ for any $\beta \in \Phi^{\bullet}_L$.
\end{construction}

\begin{proof}
The above function is constructed as follows. 
When $\zeta$ is generic, meaning $\ell=0$, all its integer powers are distinct, so the condition $\beta(x) = \zeta^{m_{\beta}}$ determines $m_{\beta} \in \Z$ uniquely. On the other hand, when $\zeta$ is a root of unity of order $\ell \geq 1$, this condition is $\ell$-periodic, so $m_{\beta}$ is unique modulo $\ell$. In both cases, we obtain a map
\begin{equation*}
    \Phi^{\bullet}_L \to \Z / \ell \Z, \qquad \beta \mapsto [m_{\beta}].
\end{equation*}
This map is additive. 
Indeed, if $\alpha(x) = \zeta^{m_{\alpha}}$ and $\beta(x) = \zeta^{m_{\beta}}$, then $(\alpha + \beta)(x) = \alpha(x)\cdot \beta(x) = \zeta^{-(m_{\alpha}+m_{\beta})}$, so $[m_{\alpha + \beta}] \equiv [m_{\alpha}] + [m_{\beta}]$ modulo $\ell$ as desired.
We now choose an additive lift
\begin{equation*}
   \Phi^{\bullet}_L \to \Z, \qquad \beta \mapsto m_{\beta}. 
\end{equation*}
When $\zeta$ is generic, there is nothing to choose. When $\zeta$ is a root of unity of $\ord(\zeta) = \ell \geq 1$, we choose arbitrarily the lifts $m_{\alpha_i} \in \Z$ of $[m_{\alpha_i}]$ for all simple roots $\alpha_i \in \Phi^{\bullet, \sr}_L$.
Since simple roots are linearly independent and generate all roots, this uniquely extends the lift to $\Phi^{\bullet}_L$ by additivity.
\end{proof}


The above choice of $\mm$ may be equivalently packaged into a cocharacter $\omega$ of the adjoint form of $L$, which we call a resonant cocharacter at $(x, \zeta)$.
\begin{construction}
Let $T^{\ad}_L$ be the maximal torus of the adjoint form $L^{\ad}$ of the resonant subgroup $L$.
We consider a cocharacter
\begin{equation*}
 \omega = \omega_{(x, \zeta)} \in X_{\bullet}(T^{\ad}_L)   
\end{equation*}
with the property that $\beta(\omega(\zeta)) = \beta(x)$, $\forall \beta \in \Phi^{\bullet}_L$ and call it a \emph{resonant cocharacter} at $(x, \zeta)$.
\end{construction}
\begin{proof}
Such a cocharacter always exists. Indeed, for each $\beta \in \Phi^{\bullet}_L$ we have $\beta(x) = \zeta^{m_{\beta}}$ for the additive function $\mm: \Phi^{\bullet}_L \to \Z$. We thus need to construct a cocharacter $\omega$ such that $\beta(\omega(x)) = \zeta^{m_{\beta}}$ for any $\beta \in \Phi^{\bullet}_L$. By additivity of $m_{\beta}$, it is enough to check this property only on simple roots $\alpha_i \in \Phi^{\bullet, \sr}_L$. These form a free basis of the root lattice $\Z\Phi^{\bullet}_L$, which is also equal to the character lattice of the adjoint form $X^{\bullet, \ad}_{L}$. Hence we have the dual basis of cocharacters $\omega_i \in X^{{\ad}}_{\bullet, L}$. Setting $\omega := \sum_i m_{\alpha_i} \cdot \omega_i$ gives the desired cocharacter.
\end{proof}





\begin{remark}\label{remark: non-uniqueness of the resonant part}
Given the choice of $\mm$, we have defined $\omega$ by the property $\langle \beta, \omega \rangle = m_{\beta}$, $\forall \beta \in \Phi^{\bullet}_L$. The ambiguity in choosing $\omega$ is the same as the ambiguity in choosing the additive function $\mm$.
\begin{enumerate}
    \item The evaluation map $\ev_{\zeta}: X^{\bullet, \ad}_{L} \to k$, $\lambda \mapsto \lambda(\zeta)$ is not injective unless $\zeta$ is generic. For generic $\zeta$, such $\omega$ is unique. For $\zeta$ root of unity of $\ord(\zeta) = \ell$, it is determined up to $\ell X_{\bullet}^{{\ad}}$.
    \item To define $\omega$, we have to pass to the adjoint form $L^{\ad}$. The constructed cocharacter $\omega$ of $T^{\ad}_{L}$ may not lift to a cocharacter of a maximal torus $T_L$ of $L$. Even if it does, such a lift of $\omega$ in $T_L$ is not unique.
\end{enumerate}
\end{remark}



\subsubsection{Parabolic and Levi subgroups of $L$.}\label{section: parabolic and levi subgrops of L}
We consider $(x, \zeta)$ as above and work inside the resonant subgroup $L=L_{[x, \zeta]}$.
%
Consider the subgroup $P_x \leq L$ spanned by the same maximal torus and root spaces 
\begin{equation}\label{equation: roots of Px}
 \Phi^{\bullet}_{P_x} := \{\beta \in \Phi^{\bullet}_L \mid \beta(x) = \zeta^m \ \text{for some} \ m \leq 0 \}
\end{equation}
For each root $\beta \in \Phi^{\bullet}_L$, at least one of $\pm \beta$ lies inside $\Phi^{\bullet}_{P_x}$. Applying this observation to simple roots of $L$, we conclude that $P_x$ contains a Borel subgroup of $L$. Consequently, $P_x \leq L$ is a parabolic subgroup.
%
By the definition of $\omega$, we further have $\beta(x) = \beta(\omega(\zeta)) = \zeta^{\langle \beta, \omega \rangle}$. Therefore, $\Phi^{\bullet}_{P_x} = \{\beta \in \Phi^{\bullet}_L \mid \langle \beta, \omega \rangle \leq 0 \}$. These are precisely the roots of the opposite parabolic subgroup $P_{\omega} \leq L$, hence $P_x = P_{\omega}$.


Moreover, for any $\beta \in \Phi^{\bullet}_L$, both $\pm \beta$ lie in $\Phi^{\bullet}_{P_x}$ if and only if $\beta(x) = 1$. This condition cuts out the root system $ \Phi^{\bullet}_{x} = \{\beta \in \Phi^{\bullet}_L \mid \beta(x) = 1 \}$ of the stabilizer subgroup $L_x$. We deduce that $L_x$ is the Levi factor of $P_x$. As in the above paragraph, this may be rewritten as $\Phi^{\bullet}_{x} = \{\beta \in \Phi^{\bullet}_L \mid \langle \beta, \omega \rangle = 0 \}$, which is the root subsystem of $L_{\omega} \leq L$.

Altogether, we have proved the following.
\begin{lemma}\label{lemma: parabolic and levi subgrops of L}
The stabilizer $L_x \leq L$ is a Levi subgroup of $L$. It is the Levi factors of the parabolic subgroup $P_x \leq L$. Moreover, $L_x = L_{\omega}$ and $P_x = P_{\omega}$ are the Levi and opposite parabolic associated to the resonant cocharacter $\omega \in X_{\bullet, L}^{{\ad}}$. 
\end{lemma}

\begin{remark}\label{remark: parabolic and levi subgrops of L}
Let $y = x \cdot \lambda(\zeta)$ for some $\lambda \in \Z\Phi^{\bullet}_G$ be an element in the $\zeta$-resonance class of $x$. Then the resonant cocharacter of $y$ is given by $\omega + \lambda^{\ad}$, where $\lambda^{\ad}$ is the projection of $\lambda$ to the cocharacter lattice of $L^{\ad}$ along $X_{\bullet, G} = X_{\bullet, L} \twoheadrightarrow X_{\bullet, L}^{\ad}$. Since $L = L_{[x, \zeta]} = L_{[y, \zeta]}$ depends only on the resonance class of $x$, Lemma \ref{lemma: parabolic and levi subgrops of L} follows for $y$.

Concretely, we deduce $L_y \leq L$ is a Levi subgroup of $L$. It is the Levi factor of the parabolic subgroup $P_{y}$. Moreover, $L_y = L_{\omega + \lambda^{\ad}}$ and $P_{\omega + \lambda^{\ad}}$ are the Levi and opposite parabolic associated to the resonant cocharacter $\omega + \lambda^{\ad} \in X_{\bullet, L}^{{\ad}}$.
\end{remark}



\subsubsection{Independence on Weyl group actions.}\label{subsubsection: independence on W-action}
Fix the point $\zeta \in \Gm^{\rot}$. Any closed point $x \in T$ projects to a closed point $\overline{x} = x \cdot W$ of $T\sslash W$. 


The resonant subgroup $L_{[x, \zeta]}$ and its Weyl group $W_{[x, \zeta]}$ depend, up to canonical isomorphism given by conjugation, only on $\overline{x}$. 
%
Since they also depend only on the $\zeta$-resonance class of $x$ by definition, they are up to canonical isomorphism attached to the orbit of the $\zeta$-integral action of the affine Weyl group $W_{\aff} \acts T$, $(w, \lambda) \cdot x = \lambda(\zeta) \cdot w(x)$. Also see \S \ref{subsubsection: resonance} and \S \ref{section: actions I} below.


Similarly, the stabilizing Weyl group $W_x$, the stabilizing subgroup $L_x$ and the associated parabolic subgroup $P_x$ of $L_{[x, \zeta]}$ depend, up to canonical isomorphism given by conjugation, only on $\overline{x}$.

In particular, we have shown the following.
\begin{claim}
The pair $P_x \leq L_{[x, \zeta]}$ depends, up to isomorphism, only on $\overline{x} \in T \sslash W$ and $\zeta \in \Gm^{\rot}$.   
\end{claim}




\subsubsection{Vocabulary to Humphrey's notation.}
When $G$ is of adjoint type and $\zeta$ is generic, the above constructions appear in the study of the classical category $\O$, see \cite{Hum08, Jan79}. More precisely, \cite{Hum08} works (under Langlands dual notation) with the BGG category $\O$ of the semisimple Lie algebra $\widecheck{\g}$. The genericity of $\zeta$ corresponds to the classical version of the category $\O$ (as opposed to its quantum version at a root of unity). In this setup, the (co)character $\omega$ is unique (as we saw in Remark \ref{remark: non-uniqueness of the resonant part}). For convenience, we provide a short vocabulary to the notation from \cite{Hum08} in this special case -- compare:
    \begin{itemize}
        \item  the definition of $W_L = W_{[x, \zeta]}$ to $W_{[c]}$ of \cite[\S 3.4]{Hum08},
        \item the definition of $\omega = \omega_{[x, \zeta]}$ to the definition of $c^{\natural}$ of \cite[\S 7.4]{Hum08},
        \item the definition of $W_x$ to the definition of $W^{\circ}_c$ of \cite[\S 7.4]{Hum08}.
    \end{itemize}






\subsection{Centralizers in loop groups}\label{subsection: centralizers in loop groups}


\subsubsection{Centralizers in loop groups.}\label{section: centralizers}

Let $(x, \zeta) \in T \times \Gm^{\rot}$. Consider its centralizer $\Cent_{\lo G}(x, \zeta) \leq \lo G$ in the loop group $\lo G$. It might be not connected; we denote the identity component by
\begin{equation}
J := J_{(x, \zeta)} := \Cent^{\circ}_{\lo G}(x, \zeta) \leq \lo G.
\end{equation}
%
Furthermore, consider the positive part the connected centralizer $J$, given by
\begin{equation}
J^+ 
:= J \cap \lop G \leq \lop G   
\end{equation}
Equivalently, $J^+ = \Stab_{J}(t^{0})$ is the stabilizer inside $J$ of the base point $t^{0} \in \Gr_G$. 

%
More generally, for any cocharacter $\lambda \in X_{\bullet}$, we consider
\begin{equation}\label{equation: positive stabilizer}
    J^{+}_{\lambda} := J \cap t^{\lambda} \cdot \lop G \cdot t^{-\lambda},
\end{equation}
the stabilizer $\Stab_{J}(t^{\lambda})$ inside $J$ of the point $t^{\lambda} \in \Gr_G$.

\begin{remark}
It will become apparent later -- see Lemmata \ref{lemma: positive stabilizers as parabolics -- generic zeta} and \ref{lemma: zeta root of unity -- positive stabilizers} -- that all $J^{+}_{\lambda}$ are already connected. For example, this a posteriori means $J^{+}$ is given by the connected centralizer $J^+ = \Cent^{\circ}_{\lop G}(x, \zeta)$ of $(x, \zeta)$ inside the positive loop group.   
\end{remark}


\subsubsection{Root system of $J$.}
We can conveniently describe the root system of $J$ as follows.
\begin{lemma}\label{lemma: root spaces of J}
Let $(x, \zeta) \in T \times \Gm^{\rot}$ arbitrary and denote $\ell= \ord(\zeta) \in \Z_{\geq 0}$. Then the root subsystem $\Phi^{\bullet}_J \subseteq \Phi^{\bullet}_{G, \aff}$ of $J = J_{(x, \zeta)}$ has the form    
\begin{align*}
\Phi^{\bullet}_J 
&= \{ \beta + m\hbar \in \Phi^{\bullet}_{G, \aff} \mid \beta(x) \cdot \zeta^{m} = 1 \} \\
&= \{\beta - (m_{\beta} + \ell \Z)\hbar \in \Phi^{\bullet}_{G, \aff} \mid \beta \in \Phi^{\bullet}_{L}\}.
\end{align*}
for the additive function $\mm: \Phi^{\bullet}_L \to \Z$, $\beta \mapsto m_{\beta}$ from Notation \ref{notation: additive function m}.
\end{lemma}

\begin{proof} 
The centralizer $J$ is determined by the closed root subsystem $\Phi^{\bullet}_J$ of the affine real root system $\Phi^{\bullet}_{G, \aff}$ of $G$ by Lemma \ref{lemma: centralizers give closed subroot systems}. Since the adjoint action of $(x, \zeta)$ on a root space $U_{\beta + m\hbar}$ is given by $\beta(x)\cdot \zeta^m$, this closed root system is given by $\Phi^{\bullet}_J = \{ \beta + m\hbar \in \Phi^{\bullet}_{G, \aff} \mid \beta(x) \cdot \zeta^{m} = 1 \}$ as in Lemma \ref{proposition: affine centralizer as regular subgroup}. The first equality is thus clear.

For the second equality, consider the resonant subgroup $L=L_{[x, \zeta]} \leq G$ with root system $\Phi^{\bullet}_L \subseteq \Phi^{\bullet}_G$. Note that $\Phi^{\bullet}_{L}$ is given by the image of $\Phi^{\bullet}_J$ under the projection \eqref{equation: projection of root systems}. Hence any root $\beta +m\hbar \in \Phi^{\bullet}_J$ has $\beta\in \Phi^{\bullet}_L$.
By Construction \ref{notation: additive function m} of the additive function $\mm$, the subset $m_{\beta} + \ell \Z \subseteq \Z$ contains precisely those integers $m$ satisfying $\beta(x) = \zeta^m$. The second equality follows.
\end{proof}

\begin{notation}\label{notation: WJ}
We denote by $W_J := W_{\Phi^{\bullet}_J} \leq W_{G, \aff}$ the reflection subgroup generated by the affine reflections along roots from $\Phi^{\bullet}_J \leq \Phi^{\bullet}_{G, \aff}$.    
\end{notation}


\begin{remark}\label{remark: root system of J project on root system of L}
Note that the resonant root system $\Phi^{\bullet}_{[x, \zeta]} = \Phi^{\bullet}_L$ from \eqref{equation: resonant root system} is the finite root subsystem given by the projection of the affine root system $\Phi^{\bullet}_J$. In terms of \eqref{equation: projection of root systems}, this reads as
\begin{equation*}
   \Phi^{\bullet}_{[x, \zeta]} = \overline{\Phi}^{\bullet}_J.
\end{equation*}
\end{remark}

\begin{remark}
Since $J$ is connected, it is determined by the root subsystem $\Phi^{\bullet}_J$. Also note that $J \leq \lo L^{\sc}$. Indeed, it is connected and contains only root spaces from $\Phi^{\bullet}_{L, \aff}$.    
\end{remark}




\subsubsection{Slanting.}\label{section: slanting}
We can get rid of the integers $m_{\beta}$ by a suitable automorphism of $\lo L$, which we call \emph{slanting}. Geometrically, slanting is given by conjugation with $t^{\omega}$, where $\omega$ is the resonant cocharacter. 

\begin{discussion}\label{discussion: slanting}
The adjoint action of a reductive group on itself factors through its adjoint form. Since $L^{\ad} \acts L$ by conjugation, we have an induced action $\lo L^{\ad} \acts \lo L$. For any cocharacter $\lambda \in X^{{\ad}}_{L, \bullet}$, conjugation by $t^{\lambda}$ gives an automorphism of $\lo L$. We call this map \emph{slanting by $\lambda$} and denote it
\begin{equation}
    \slant_{\lambda}: \lo L \to \lo L, \qquad g \mapsto t^{\lambda}\cdot g \cdot t^{-\lambda}
\end{equation}
By the same argument, slanting restricts to an automorphism of the identity component $\slant_{\lambda}: \lo L^{\sc} \to \lo L^{\sc}$. On the level of its root system, slanting is induced by
\begin{equation}\label{equation: slanting on root data}
    \slant_{\lambda}: \Phi^{\bullet}_{L, \aff} \to \Phi^{\bullet}_{L, \aff}, 
    \qquad
    \beta + m\hbar \mapsto \beta + (m + \langle \beta, \lambda \rangle)\hbar.
\end{equation}
Clearly, $\slant_{\lambda}$ and $\slant_{-\lambda}$ are inverse to each other.
%
Using this notation, any subgroup $H \leq \lo L$ can be slanted to an isomorphic copy $\slant_{\lambda}(H) := t^{\lambda} \cdot H \cdot t^{-\lambda} \leq \lo L$.

Furthermore, slanting on root systems \eqref{equation: slanting on root data} induced a map
\begin{equation}\label{equation: slanting on affine weyl groups}
    \slant_{\lambda}: W_{L, \aff} \to W_{L, \aff}.
\end{equation}
Any subgroup of $W_{L, \aff}$ may be again slanted by the automorphism $\slant_{\omega}$. Under this automorphism, the natural action $W_{L, \aff} \acts X_{\bullet}$ identifies with the $\lambda$-shifted action $(\slant_{\lambda}(W_{L, \aff}, \bullet_{\lambda}) \acts X_{\bullet}$.
Indeed, this is clear from the defining formulas \S \ref{section: shifted actions}, \S \ref{section: shifted affine weyl group actions} and \eqref{equation: slanting on root data}.
\end{discussion}

 



In particular, we may slant by the resonant cocharacter $\omega$. 
Slanting $J \leq \lo L^{\sc}$, we obtain the embedding $\slant_{\omega}(J) = t^{\omega} \cdot J \cdot t^{-\omega} \leq \lo L^{\sc}$.
As we saw, this embedding can be defined purely on the level of associated root systems.
%
\begin{lemma}\label{discussion: slanting J in general}
The closed root subsystem of $\slant_{\omega}(J)$ is given by
\begin{equation*}
    \Phi^{\bullet}_{\slant_{\omega}(J)} = \{ \beta + \ell \Z \mid \beta \in \Phi^{\bullet}_L \}.
\end{equation*} 
\end{lemma}
\begin{proof}
Since conjugation by $t^{\omega}$ sends the root space $U_{\beta + m\hbar}$ to $U_{\beta + (m + \langle \omega, \beta \rangle \hbar )} = U_{\beta + (m + m_{\beta}) \hbar}$ by Lemma \ref{lemma: conjugation of affine root spaces}, the slanting map on root systems identifies
\begin{align*}
   \Phi^{\bullet}_J &\xrightarrow{=} \Phi^{\bullet}_{\slant_{\omega}(J)} \\
   \beta + m\hbar &\mapsto \beta + (m + m_{\beta})\hbar.
\end{align*}
The result then follows from the description of $\Phi^{\bullet}_J$ in Lemma \ref{lemma: root spaces of J}.
\end{proof}

The slanted subgroup $t^{\omega} \cdot J \cdot t^{-\omega}$ is simpler to manipulate -- as we will record in Lemmata \ref{lemma: centralizer as resonant subgroup -- generic zeta}, \ref{lemma: zeta root of unity -- computation of centralizer}, it is given by the reductive group $L$ resp. the identity component of its dilated loop group, embedded in the naive way. Vice versa, the original centralizer $J$ is then given by slanting these naive embeddings via $\slant_{-\omega}$.


\begin{remark}\label{remark: slanting on weyl groups general case}
The slanting map on root data $\slant_{\omega}: \Phi^{\bullet}_J \xrightarrow{=} \Phi^{\bullet}_{\slant_{\omega}(J)}$ induced a map on the corresponding reflection subgroups
\begin{equation*}
    \slant_{\omega}: W_J \xrightarrow{=} W_{\slant_{\omega}(J)}
\end{equation*}
of $W_{L, \aff}$. By Discussion \ref{discussion: slanting}, $\slant_{\omega}$ intertwines the naive action and the $\omega$-shifted action on $X_{\bullet}$.
\end{remark}




\subsection{Weyl group stabilizers for simply connected \texorpdfstring{$G$}{G}}\label{subsection: weyl group stabilizers}
Let $k$ be algebraically closed of characteristic zero.
Assuming $G = G^{\sc}$ is simply connected, we discuss stabilizers of finite and affine Weyl group actions on its maximal torus $T = T^{\sc}$. The simply connected assumption is needed for technical reasons -- even in the reductive situation, one has to consider stabilizers for $T^{\sc}$.



\subsubsection{Affine Weyl group action on extended torus.}\label{section: actions I}
We have the defining action of the finite Weyl group $W \acts T$. It lifts to a natural action of the affine Weyl group on the extended torus
\begin{align}\label{equation: affine weyl group action I}
\begin{split}
W_{\aff} &\acts \Gm^{\rot} \times T \\
(\lambda, w) \cdot (\zeta, x) &= (\zeta, \lambda(\zeta) \cdot w(x)).
\end{split}
\end{align}
It is useful to view this as an algebraic family over $\Gm^{\rot}$ of actions of $W_{\aff} = \Z\Phi^{\bullet} \rtimes \Gm^{\rot}$ on the torus $T$. Namely, for fixed $\zeta \in \Gm^{\rot}$, we have the \emph{$\zeta$-integral action}
\begin{equation}\label{equation: zeta integral action of affine Weyl group}
W_{\aff} \acts T, \qquad (\lambda, w) \cdot x = \lambda(\zeta) \cdot w(x).    
\end{equation}


\subsubsection{Weyl group stabilizers as reflection subgroups.}
Consider a closed point $x \in T=T^{\sc}$.
\begin{lemma}\label{lemma: stabilizer in finite weyl group}
The stabilizer of $x$ under the action of $W \acts T$ is given by
\begin{equation*}
    \Stab_W(x) = W_x.
\end{equation*}   
\end{lemma}
\begin{proof}
The stabilizer $\Stab_W(x) \leq W$ is a reflection subgroup by \S \ref{section: stabilizers in weyl groups as reflection subgroups}, Proposition \ref{proposition: steinberg fixed point theorem}.
Hence we only have to describe which reflections $s_{\beta}$, $\beta \in \Phi^{\bullet}_G$ lie inside it. Since $s_{\beta}(x) = x \cdot \widecheck{\beta}(\beta(x))^{-1}$, we see that $s_{\beta}$ fixes $x$ if and only if $\beta(x)=1$. Consequently, $\Stab_W(x)$ is the reflection subgroup of the root subsystem $\{ \beta \in \Phi^{\bullet}_G \mid \beta(x)=1 \}$. By definition, this is the root subsystem $\Phi^{\bullet}_x$ defining $W_x \leq W$. Also see Proposition \ref{proposition: pseudo-levi subgroups}, (iii).
\end{proof}

Taking a closed point $(x, \zeta) \in T \times \Gm^{\rot} = T^{\sc} \times \Gm^{\rot}$, we have the following affine analogue. To state it, recall from Notation \ref{notation: WJ} that $W_J = W_{\Phi^{\bullet}_J}$ stands for the reflection subgroup of $W_{\aff}$ associated to the root subsystem $\Phi^{\bullet}_J$ of $J$. This root subsystem was described in Lemma \ref{lemma: root spaces of J}.

\begin{lemma}\label{lemma: affine stabilizer in weyl group}
The stabilizer of $(x, \zeta)$ under the action $W_{\aff} \acts T \times \Gm^{\rot}$ is given by
\begin{equation*}
    \Stab_{W_{\aff}}(x, \zeta) = W_J.
\end{equation*}
In other words, the stabilizer of $x \in T$ under the $\zeta$-integral action of $W_{\aff}$ is given by $W_J$.
\end{lemma}
\begin{proof}
Fix $\zeta \in \Gm^{\rot}$. Then $\Stab_{W_{\aff}}(x, \zeta)$ is given by the stabilizer $\Stab_{W_{\aff}}(x)$ under the $\zeta$-integral action \eqref{equation: zeta integral action of affine Weyl group}. 
As in the finite case, the stabilizer $\Stab_{W_{\aff}}(x)$ is a reflection subgroup by Proposition \ref{proposition: steinberg fixed point theorem}. 
%
Hence it is enough to determine which affine reflections $s_{\beta + m\hbar} \in \Phi^{\bullet}_{G, \aff}$ lie inside it. Since $s_{\beta + m\hbar}(x) = x \cdot \widecheck{\beta}(\beta(x) \cdot \zeta^m)^{-1}$, the reflection $s_{\beta + m\hbar}$ fixes $x$ if and only if $\beta(x) \cdot \zeta^m =1$. Hence $\Stab_{W_{\aff}}(x)$ is the reflection subgroup associated to the root system $\{\beta + m\hbar \in \Phi^{\bullet}_{G, \aff} \mid \beta(x) \cdot \zeta^m =1 \}$. Since this is the root system $\Phi^{\bullet}_J$ of $J$, the result follows.
\end{proof}

As a corollary, we deduce that the orbit of $x \in T$ under the $\zeta$-integral action $W_{\aff} \acts T$ is given by $W_{\aff} / W_J$. These orbits will be discussed further in \S \ref{section: some useful notation} and prominently appear as the labeling set for connected componets of fixed points on $\Gr_G$.


\subsubsection{Image of stabilizer in finite Weyl group.}

We may easily describe the image of the stabilizer under the projection to the finite Weyl group.

\begin{lemma}
Let $(x, \zeta) \in T \times \Gm^{\rot}$.
Under the projection map $W_{\aff} \twoheadrightarrow W$, the image of $\Stab_{W_{\aff}}(x, \zeta)$ is given by the resonant Weyl group $W_{[x, \zeta]} \leq W$.
\end{lemma}
\begin{proof}
The projection in question is induced by the projection of root systems $\Phi^{\bullet}_{G, \aff} \twoheadrightarrow \Phi^{\bullet}_{G}$. By Lemma \ref{lemma: affine stabilizer in weyl group}, $\Stab_{W_{\aff}}(x, \zeta) = W_{J}$ is the reflection subgroup associated to the subsystem $\Phi^{\bullet}_J$, whose projection is precisely the root subsystem $\Phi^{\bullet}_{[x, \zeta]} $ of $W_{[x, \zeta]}$ by Remark \ref{remark: root system of J project on root system of L}. The result follows.
\end{proof}

If $\zeta \in \Gm - \mu_{\infty}$ is generic, the above stabilizer $\Stab_{W_{\aff}}(x, \zeta) = W_J$ thus fits in the following diagram whose rows are short exact sequences.
\begin{equation}\label{equation: stabilizer exact sequences for zeta generic}
    \begin{tikzcd}
        1 \arrow[r] & \Z\Phi_{\bullet} \arrow[r] & W_{\aff} \arrow[r] & W \arrow[r] & 1 \\
        1 \arrow[r] & 1 \arrow[r] \arrow[u, hook] & \Stab_{W_{\aff}}(x, \zeta) \arrow[r , "="] \arrow[u, hook] & W_{[x, \zeta]} \arrow[r] \arrow[u, hook] & 1
    \end{tikzcd}
\end{equation}

On the other hand, assume now $\zeta \in \mu_{\infty}$ is primitive root of unity of $\ord(\zeta) = \ell$. Then $\Stab_{W_{\aff}}(x, \zeta) = W_J$ fits in the diagram of horizontal short exact sequences
\begin{equation}\label{equation: stabilizer exact sequences for zeta root of unity}
    \begin{tikzcd}
        1 \arrow[r] & \Z\Phi_{\bullet} \arrow[r] & W_{\aff} \arrow[r] & W \arrow[r] & 1 \\
        1 \arrow[r]& \ell \Z\Phi_{\bullet} \arrow[r] \arrow[u, hook] & \Stab_{W_{\aff}}(x, \zeta) \arrow[r] \arrow[u, hook] & W_{[x, \zeta]} \arrow[r] \arrow[u, hook] & 1 \\
        1 \arrow[r]  & \ell \Z\Phi_{\bullet} \arrow[r] \arrow[u, hook] & \ell \Z\Phi_{\bullet} \rtimes W_{x} \arrow[r] \arrow[u, hook] & W_{x} \arrow[r] \arrow[u, hook] & 1 
    \end{tikzcd}
\end{equation}
Here, the bottom row is given by the separate stabilizers of $x$ inside $\Z\Phi_{\bullet}$ and $W$. Note that the actual stabilizer in the middle row is bigger that their semidirect product.



To wrap up, if $G$ is simply connected, the reflection subgroups \eqref{equation: resonant weyl group} are equivalently given by
\begin{equation}\label{eq: weyl group and resonant weyl group}
   W_x = \{w \in W \mid wx = x\} 
   \qquad \text{and} \qquad  
   W_{[x, \zeta]} = \{w \in W \mid wx \sim_{\zeta} x\}.
\end{equation}
In words, $W_x$ is the stabilizer of $x \in T$, while $W_{[x, \zeta]}$ is the stabilizer of its resonance class $[x, \zeta]$.     







\subsection{Cocharacter graphs for simply connected \texorpdfstring{$G$}{G}}\label{section: cocharacter graphs}
The above combinatorics is conveniently encoded by the \textit{cocharacter graphs} $\overline{\Gamma} \subseteq T \times \Gm^{\rot} \times T \sslash W$ introduced in this section. 

To simplify the exposition, we will again assume $G =G^{\sc}$ to be simply connected throughout \S \ref{section: cocharacter graphs}. With extra bookkeeping, the same discussion works for any reductive group; the necessary modifications are described afterwards in \S \ref{section: cocharacter graphs for reductive groups}.


\subsubsection{Cocharacter graphs for simply connected groups.}\label{subsubsection: cocharacter graphs for simply connected groups}
Assume $G$ is simply connected. Then $\Z\Phi^{\bullet}_G = X_{\bullet}$ and $W_{\aff} = W_{\ext}$. Recall the action \eqref{equation: affine weyl group action I}. For any element $(\lambda, w) \in W_{\aff}$ with $\lambda \in X_{\bullet}$ and $w \in W$, consider the map
\begin{align*}
    (\lambda, w): T \times \Gm^{\rot} &\to T \\
    (x, \zeta) &\mapsto  \lambda(\zeta) \cdot w(x).
\end{align*}
and the closed subscheme $\Gamma_{\lambda, w} \hookrightarrow T \times \Gm^{\rot} \times T$ given by its graph
\begin{align*}
    \Gamma_{\lambda, w} 
    := &\{ (x, \zeta, \lambda(\zeta) \cdot w(x)) \in T \times \Gm^{\rot} \times T \} \\
    = &\{ (x, \zeta, y) \in T \times \Gm^{\rot} \times T \mid y = \lambda(\zeta) \cdot w(x) \}.
\end{align*}
Taking the union of these closed subschemes over varying $\lambda$, we obtain the closed ind-subscheme
\begin{equation*}
    \Gamma := \{ (x, \zeta, y) \in T \times \Gm^{\rot} \times T \mid \exists \lambda \in X_{\bullet}, \ y \cdot w(x)^{-1} = \lambda(\zeta) \} \hookrightarrow T \times \Gm^{\rot} \times T.
\end{equation*}
The ind-scheme $\Gamma$ is reduced. Its irreducible componets $\Gamma_{\lambda, w}$ are copies of $T \times \Gm^{\rot}$ labeled by $W_{\aff}$, intersecting in various ways. In particular, we have the distinguished component $\Gamma_{1, 1}$, parametrizing the points $\{ (x, \zeta, x) \in T \times \Gm^{\rot} \times T\}$.


\subsubsection{Action of affine Weyl group.}\label{section: actions}
Consider the action of $W_{\aff}$ on $T \times \Gm^{\rot} \times T$ given by the trivial action on the first copy of $T$ and the natural action \eqref{equation: affine weyl group action I} on $\Gm^{\rot} \times T$, namely
\begin{equation*}
   W_{\aff} \acts T \times \Gm^{\rot} \times T, \qquad (\lambda, w) \cdot (x, \zeta, y) = (x, \zeta, \lambda(\zeta) \cdot w(y)). 
\end{equation*}
This action restricts to an action
\begin{equation}\label{equation: affine weyl group action on cocharacter graphs}
W_{\aff} \acts \Gamma    
\end{equation}
on the closed sub-ind scheme $\Gamma \hookrightarrow T \times \Gm^{\rot} \times T$. It acts freely and transitively on the set of irreducible components $(\Gamma_{\lambda, w} \mid (\lambda, w) \in W_{\aff})$ of $\Gamma$.

\subsubsection{Variants.}\label{section: variants}
Given any pair of subgroups $W_1, W_2 \leq W$, consider their action on $T \times T$. The ind-subscheme $\Gamma \hookrightarrow T \times \Gm^{\rot} \times T$ is clearly $(W_1 \times W_2)$-equivariant. Taking the quotient, we get the sub-ind scheme
\begin{equation*}
    \Gamma^{T \sslash W_1,  T \sslash W_2} \hookrightarrow T\sslash W_1 \times \Gm^{\rot} \times T \sslash W_2.
\end{equation*}
We will be particularly interested in the case $W_1= 1$ and $W_2 = W$. We then write
\begin{equation*}
    \overline{\Gamma} := \Gamma^{T,  T \sslash W} \hookrightarrow T \times \Gm^{\rot} \times T \sslash W.
\end{equation*}

The action of $W_2=W$ on the right copy of $T$ permutes the irreducible components of $\Gamma$ according to the left translation on the labeling set $W \acts (X_{\bullet} \rtimes W)$. In other words, $\forall w \in W$, we have $w \cdot \Gamma_{\lambda, v} = \Gamma_{w(\lambda), wv}$.
The irreducible components of $\overline{\Gamma}$ are thus copies of $T \times \Gm^{\rot}$, labeled by the quotient set $X_{\bullet} = W \backslash (X_{\bullet} \rtimes W)$. They correspond to $W$-orbits of irreducible components of $\Gamma$. 
For each $\lambda \in X_{\bullet}$, we denote by $\overline{\Gamma}_{\lambda}$ the component corresponding to the $W$-orbit of $\Gamma_{\lambda, 1}$.   



\subsubsection{Fibers.}\label{section: fibers}
Given $\Gamma$, $\overline{\Gamma}$ and the other variants, we may consider their fibers over points in the base $T \times \Gm^{\rot}$. 
%
Given any $(x, \zeta) \in T \times \Gm^{\rot}$, we denote these fibers as
\begin{equation*}
    \Gamma_{(x, \zeta)} \hookrightarrow T 
    \qquad
    \text{and}
    \qquad
    \overline{\Gamma}_{(x, \zeta)} \hookrightarrow T\sslash W.
\end{equation*}
These are closed-ind schemes. Their underlying reduced schemes are given by, possibly infinite, union of points.
%
Similarly, we denote the fibers over some $\zeta \in \Gm^{\rot}$ by
\begin{equation*}
    \Gamma_{\zeta} \hookrightarrow  T \times T 
    \qquad
    \text{and}
    \qquad
    \overline{\Gamma}_{\zeta} \hookrightarrow T \times T\sslash W.
\end{equation*}
These are again closed sub-ind-schemes. The irreducible components of the reduces scheme underlying $\Gamma_{\zeta}$ are given by copies of $T$.




\subsubsection{Fibers of $\Gamma$ as orbits on cocharacter lattice.}\label{section: specializations of the lattice}
For any $(x, \zeta) \in T \times \Gm^{\rot}$, the action $W_{\aff} \acts \Gamma$ from \eqref{equation: affine weyl group action on cocharacter graphs} restricts to a transitive action
\begin{equation}\label{equation: transitive action of Waff}
  W_{\aff} \acts |\Gamma_{(x, \zeta)}|
\end{equation}
on the discrete set of closed points of the fiber $\Gamma_{(x, \zeta)}$. Under the embedding to $|T|$, the set $|\Gamma_{(x, \zeta)}|$ is thus given by the $W_{\aff}$-orbit of $x$ under the $\zeta$-integral action \eqref{equation: zeta integral action of affine Weyl group}. 

\begin{lemma}\label{lemma: identification of fibers of Gamma}
We have a natural identification  
\begin{equation*}
 |\Gamma_{(x, \zeta)}| = (X_{\bullet} \rtimes W) /  W_J.
\end{equation*}
\end{lemma}
\begin{proof}
The stabilizer in $W_{\aff}$ of the base point $(x, \zeta, x) \in |\Gamma_{(x, \zeta)}| \subseteq |T \times \Gm^{\rot} \times T|$ is given as
\begin{equation*}
    \Stab_{W_{\aff}}(x, \zeta, x) = W_J.
\end{equation*}
by Lemma \ref{lemma: affine stabilizer in weyl group}. Considering the transitive action \eqref{equation: transitive action of Waff}, we conclude.
\end{proof}


\begin{remark}\label{remark: colisions of components of Gamma}
In other words: the irreducible components of $\Gamma$ are copies of $T \times \Gm^{\rot}$ labeled by $W_{\aff}$, whose collisions over the point $(x, \zeta)$ are described by the quotient map
\begin{equation*}
    W_{\aff} \twoheadrightarrow W_{\aff} / W_{J}.
\end{equation*}
Similarly, the irreducible components of the reduced ind-scheme underlying $\Gamma_{\zeta}$ are copies of $T$ labeled by $W_{\aff}/ \ell X_{\bullet}$. whose collisions over the point $x \in T$ are given by
$W_{\aff}/ \ell X_{\bullet} \twoheadrightarrow  (W_{\aff}/ \ell X_{\bullet})/ W_{L}$.
\end{remark}




We may also describe the fiber $|\overline{\Gamma}_{(x, \zeta)}| \hookrightarrow |T\sslash W|$ similarly to Lemma \ref{lemma: identification of fibers of Gamma}. The group $W_{\aff}$ acts on itself by left and right translations. In particular, we get the commuting actions
\begin{equation*}
    W \acts W_{\aff} \actsr W_J
\end{equation*}
Here, the left action is induced by the given embedding of $W$ into the semidirect product defining $W_{\aff}$. The quotient by this $W$ naturally identifies as $X_{\bullet} = W\backslash W_{\aff}$. This identification is compatible with the right actions of the stabilizer $W_J \leq W_{\aff}$.

\begin{lemma}\label{lemma: fiber of cocharacter graphs as lattice quotient}
We have a natural identification    
\begin{equation*}
 |\overline{\Gamma}_{(x, \zeta)}| = X_{\bullet} /  W_J.
\end{equation*}
\end{lemma}
\begin{proof}
With the above left $W$-action, the embedding
$W_{\aff} / \Stab_{W_{\aff}}((x, \zeta, x)) = |\Gamma_{(x, \zeta)}| \hookrightarrow T$
becomes $W$-equivariant. The claim follows from Lemma \ref{lemma: identification of fibers of Gamma} by taking quotient.    
\end{proof}

We make this quotient more explicit in Lemmata \ref{lemma: labeling set for components -- zeta generic}, \ref{lemma: labeling set for components -- zeta root of unity}, depending on the genericity of $\zeta$.




\subsubsection{Labeling set for components.}\label{section: some useful notation}
\begin{notation}\label{notation: Theta}
For notational convenience, we will denote the discrete set from Lemma \ref{lemma: fiber of cocharacter graphs as lattice quotient} by
\begin{equation*}
    \Theta := |\overline{\Gamma}_{(x, \zeta)}| = X_{\bullet} /  W_J.
\end{equation*}
Note that $\Theta$ implicitly depends on $(x, \zeta)$. An element $\vartheta \in \Theta$ may be equivalently regarded as:
\begin{enumerate}
    \item[(i)] a closed point of the fiber of $\overline{\Gamma}$ over $(x, \zeta)$ (by definition);
    \item[(ii)] a closed point of $|T\sslash W|$ which is $\zeta$-resonant with $x$ modulo $W$ (under $ |\overline{\Gamma}_{(x, \zeta)}| \subseteq |T \sslash W|$);
    \item[(iii)] an element in the quotient $X_{\bullet} /  W_J$ (under the orbit description of Lemma \ref{lemma: fiber of cocharacter graphs as lattice quotient})
    \item[(iv)] a closed point of the quasi-diagonal $Z_{\ell} \subseteq T \times T \sslash W$ (as in \S \ref{section: quasi-diagonals at roots of unity} by Lemma \ref{lemma: Z is Gamma} below).
\end{enumerate} 
\end{notation}

The set $\Theta$ will play an essential role in \S \ref{section: fixed points on affine grassmannians} as the labeling set for components of fixed points on the affine Grassmannian. To each point of $\Theta$, we can associate the stabilizing subgroups as follows.

\begin{notation}\label{notation: levi and paraolic subgroups labeled by elements of Theta}
Given $x\in T$, recall from \S \ref{subsubsection: independence on W-action} that the Weyl group $W_x$, the stabilizing subgroup $L_x$ and the associated parabolic subgroup $P_x$ of $L_{[x, \zeta]}$ depend, up to canonical isomorphism, only on its image $\overline{x}$ in $T \sslash W$.
%
Interpreting $\vartheta \in \Theta$ as a closed point of $T \sslash W$ with a lift $y \in T$, it is thus consistent to denote
\begin{equation*}
    W_{\vartheta} := W_{y} , \qquad
    L_{\vartheta} := L_{y} , \qquad
    P_{\vartheta} := P_{y} 
\end{equation*}
We have $W_{\vartheta} \leq W_{[x, \zeta]}$ and $L_{\vartheta} \leq P_{\vartheta}  \leq L_{[x, \zeta]}$. The isomorphism class of the subgroups $W_{\vartheta}$, $L_{\vartheta}$, $P_{\vartheta}$ varies as $\vartheta$ goes through $\Theta$.

We would like to stress that $W_{\vartheta}$, $L_{\vartheta}$, $P_{\vartheta}$ are \textbf{not} the subgroups associated to a cocharacter $\lambda \in X_{\bullet}$ lifting $\vartheta \in X_{\bullet}/W_J$. Rather, these subgroups are associated to the point $x \cdot \lambda(\zeta)$, or alternatively to the adjoint cocharacter $\omega + \lambda^{\ad} \in X^{\ad}_{\bullet, L}$. There is an extra shift by $\omega$.
\end{notation}




\subsection{Cocharacter graphs for reductive groups}\label{section: cocharacter graphs for reductive groups}

\subsubsection{Comments.}
Assume now $G$ to be reductive. For technical reasons, the above discussion needs to be modified. Firstly, we need to replace $W_{\aff}$ by $W_{\ext} = X_{\bullet} \rtimes W$. Furthermore, we need to replace the second copy of $T$ by $T^{\sc} \times \pi_1(G)$, the product of its simply connected cover and the fundamental group of $G$. At first in \S \ref{subsubsection: cocharacter graphs for reductive groups}, \S \ref{subsubsection: variants and stabilizers for reductive G}, we will assume $G$ is semisimple and also replace the first copy of $T$ by $T^{\sc}$; we show how to descend back to $T$ in \S \ref{section: action of the fundamental group}.

\subsubsection{Cocharacter graphs for reductive groups.}\label{subsubsection: cocharacter graphs for reductive groups}
We start by exteding the definitions from \S \ref{subsubsection: cocharacter graphs for simply connected groups}.
First, note that a character $\lambda: \Gm^{\rot} \to T$ lifts to a character $\Gm^{\rot} \to T^{\sc}$ if and only if its class in the fundamental group $\overline{\lambda} \in \pi_1(G)$ is trivial. As in \S \ref{section: borovoi fundamental group}, we denote $\chi(\overline{\lambda}) \in X_{\bullet}(T)$ a chosen representative of $\overline{\lambda}$ under the surjection $X_{\bullet}(T) \twoheadrightarrow \pi_1(G)$. Consequently, $\underline{\lambda} := \lambda \cdot \chi(\overline{\lambda})^{-1}$ lies in $X_{\bullet}(T^{\sc}) \leq X_{\bullet}(T)$.
%
For any $(\lambda, w) \in W_{\ext}$ with $\lambda \in X_{\bullet}(T)$ and $w \in W$, we thus have the map
\begin{align*}
    (\lambda, w): T^{\sc} \times \Gm^{\rot} &\to T^{\sc} \times \pi_1(G) \\
    (x, \zeta) &\mapsto (\underline{\lambda}(\zeta) \cdot w(x), \overline{\lambda}).
\end{align*}
In words, by taking $\underline{\lambda}$, we consider the cocharacter $\lambda$ as a map into the copy of $T^{\sc}$ labeled by $\overline{\lambda} \in \pi_1(G)$. We then consider the closed subscheme given by its graph
\begin{equation*}
    \Gamma_{\lambda, w} := \{ (x, \zeta, y, \sigma) \in T^{\sc} \times \Gm^{\rot} \times T^{\sc} \times \pi_1(G) \mid y = \underline{\lambda}(q) \cdot w(x), \ \sigma = \overline{\lambda} \} \hookrightarrow T^{\sc} \times \Gm^{\rot} \times T^{\sc} \times \pi_1(G).
\end{equation*}
%
Taking the union of these closed subschemes over all $\lambda \in X_{\bullet}$, $w \in W$, we obtain the closed ind-subscheme
\begin{equation*}
    \Gamma \! := \! \{ (x, \zeta, y, \sigma) \in T^{\sc} \times \Gm^{\rot} \times T^{\sc} \times \pi_1(G) \mid \exists \lambda \in X_{\bullet}, \ y = \underline{\lambda}(\zeta)\cdot w(x), \ \sigma = \overline{\lambda} \} \! \hookrightarrow \! T^{\sc} \times \Gm^{\rot} \times T^{\sc} \times \pi_1(G).
\end{equation*}
In words, the ind-subscheme $\Gamma$ now lives inside a disjoint union of copies of $T^{\sc} \times \Gm^{\rot} \times T^{\sc}$ labeled by the discrete set $\pi_1(G)$. For each $\sigma \in \pi_1(G)$, we denote $\Gamma^{\sigma} \hookrightarrow T^{\sc} \times \Gm^{\rot} \times T^{\sc} \times \{\sigma\}$ the corresponding component. While all these copies are abstractly isomorphic via the choice of lifts $\chi(\sigma)$, these identifications are not canonical. 


The irreducible components $\Gamma_{\lambda, w}$ of $\Gamma$ are copies of $T^{\sc} \times \Gm^{\rot}$ labeled by $W_{\ext}$, intersecting in various ways. In particular, we have the distinguished component $\Gamma_{1, 1}$, parametrizing the quadruples $\{ (x, \zeta, x, 1) \in T^{\sc} \times \Gm^{\rot} \times T^{\sc} \times \pi_1(G)\}$.

There is again a natural action 
\begin{equation*}
W_{\ext} = X_{\bullet} \rtimes W \acts \Gamma, \qquad  (\lambda, w) \cdot (x, \zeta, y, \sigma) = (x, \zeta, \underline{\lambda}(\zeta) \cdot w(y), \overline{\lambda} \cdot \sigma ).   
\end{equation*}
In particular, it acts freely and transitively on the set of irreducible components of $\Gamma$. 



\subsubsection{Variants and stabilizers.}\label{subsubsection: variants and stabilizers for reductive G}
We repeat the discussion \S \ref{section: actions}, \S \ref{section: variants}, \S \ref{section: fibers}: after taking the quotient, we get the variants $\Gamma^{T^{\sc} \sslash W_1, T^{\sc} \sslash W_2} \hookrightarrow T^{\sc} \sslash W_1 \times \Gm \times T^{\sc} \sslash W_2 \times \pi_1(G)$. We in particular denote $\overline{\Gamma} := \Gamma^{T^{\sc}, T^{\sc} \sslash W} \hookrightarrow T^{\sc} \times \Gm^{\rot} \times T^{\sc} \sslash W \times \pi_1(G)$. We can further take fibers, consider stabilizers, etc.

Replacing all occurrences of $W_{\aff}$ by $W_{\ext}$, the results of \S \ref{section: specializations of the lattice} go through with the following extra argument. We again consider the stabilizer of the base point $\Stab_{W_{\ext}}((x, 1, x, 1))$. The extended affine Weyl group $W_{\ext}$ acts on the components of $T^{\sc} \times \Gm^{\rot} \times T^{\sc} \times \pi_1(G)$ labeled by $\pi_1(G)$ through the surjection $W_{\ext} \twoheadrightarrow W_{\ext} / W_{\aff} = \pi_1(G)$. After fixing $\sigma = 1$, we are thus left with the previous action of $W_{\aff} = \Z\Phi^{\bullet}_G \rtimes W$ (which is the extended affine Weyl group of $G^{\sc}$).
%
In particular, Lemma \ref{lemma: fiber of cocharacter graphs as lattice quotient} holds; we again write
\begin{equation*}
 \Theta = |\overline{\Gamma}_{(x, \zeta)}| = X_{\bullet} /  W_J
\end{equation*}
and use the further notation from \S \ref{section: some useful notation}, with respect to the points of the (appropriate copy of) $T^{\sc}$.




\subsubsection{Action of the fundamental group.}\label{section: action of the fundamental group}
The fundamental group $\pi_1(G)$ acts on $T^{\sc}$ through the central deck transformations
\begin{equation*}
    \act: \pi_1(G) \to \pi_1(G^{\ad}) = Z(G^{\sc}) \leq T^{\sc} \acts T^{\sc}.
\end{equation*}
Consequently, $\pi_1(G) \acts T^{\sc} \times \Gm^{\rot} \times T^{\sc} \times \pi_1(G)$, using the diagonal deck transformation $\act \times \act$ on $T^{\sc} \times T^{\sc}$ and the trivial action on $\Gm^{\rot} \times \pi_1(G)$. 

This action preserves each of the closed subschemes $\Gamma_{\lambda, w}$ for $(\lambda, w) \in W_{\ext}$. Indeed, $\Gamma_{\lambda, w}$ parametrizes $(x, \zeta, y, \sigma)$ satisfying $y = \lambda(\zeta) \cdot w(x)$ and $\sigma = \overline{\lambda}$. An element $\tau \in \pi_1(G)$ acts by $\tau \cdot (x, \zeta, y, \sigma) = (\tau x, \zeta, \tau y, \sigma)$. The latter condition is preserved trivially; for the former one note that
\begin{equation*}
    y = \lambda(\zeta) \cdot w(x) 
    \iff \tau \cdot y = \tau \cdot \lambda(\zeta) \cdot w(x) 
    \iff  \tau \cdot y = \lambda(\zeta) \cdot w(\tau \cdot x)
\end{equation*}
since the action of $W$ is trivial on $Z(G^{\sc})$ and compatible with the product on $T^{\sc}$. 

Altogether, $\pi_1(G)$ acts on $\Gamma \hookrightarrow T^{\sc} \times T^{\sc} \times \Gm^{\rot} \times \pi_1(G)$. This action is clearly free. We obtain the quotient
\begin{equation}\label{equation: descended Gamma eq 1}
    \Gamma / \pi_1(G) \hookrightarrow \tfrac{T^{\sc} \times T^{\sc}}{\pi_1(G)} \times \Gm^{\rot} \times \pi_1(G). 
\end{equation}

\begin{remark}\label{remark: identification for Gamma in semisimple case}
We may further use the following identification. Consider the group isomorphism $T^{\sc} \times T^{\sc} \xrightarrow{\cong} T^{\sc} \times T^{\sc}$ given by $(x, y) \mapsto (x, x^{-1}y)$. Under this isomorphism, the diagonal action $\act \times \act$ of $\pi_1(G)$ identifies with its action on the first factor $\act \times \id$. Taking the quotient identifies $\tfrac{T^{\sc} \times T^{\sc}}{\pi_1(G)} = \tfrac{T^{\sc}}{\pi_1(G)} \times T^{\sc} = T^{\der} \times T^{\sc}$. We rewrite \eqref{equation: descended Gamma eq 1} as
\begin{equation}\label{equation: descended Gamma eq 2}
    \Gamma^{T^{\der}, T^{\sc}} \hookrightarrow T^{\der} \times \Gm^{\rot} \times T^{\sc} \times \pi_1(G). 
\end{equation}
where the sub-ind scheme $\Gamma^{T^{\der}, T^{\sc}}$ is given by transporting $\Gamma$ through the above identifications and taking the quotient by $\pi_1(G)$.
\end{remark}

 

\begin{remark}\label{remark: identification for Gamma in general reductive case}
For semisimple $G$, we have $T=T^{\der}$ and Remark \ref{remark: identification for Gamma in semisimple case} already describes the desired cocharacter graphs.

For reductive $G$, we consider \eqref{equation: descended Gamma eq 2} for $G^{\ad}$, giving $\Gamma^{T^{\ad}, T^{\sc}} \hookrightarrow T^{\ad} \times \Gm^{\rot} \times T^{\sc} \times \pi_1(G)$. By base change along the map $T \to T^{\ad}$ to the first factor we obtain
\begin{equation}\label{equation: descended Gamma eq 3}
    \Gamma^{T, T^{\sc}} \hookrightarrow T \times \Gm^{\rot} \times T^{\sc} \times \pi_1(G),
\end{equation}
the appropriate generalization of the cocharacter graphs for the reductive group $G$.
In particular, we obtain the variant $\overline{\Gamma}^{T, T^{\sc}\sslash W} \hookrightarrow T \times \Gm^{\rot} \times T^{\sc}\slash W \times \pi_1(G)$. 
\end{remark}

